\documentclass[UTF8,a4paper,11pt]{ctexart}

\usepackage{amsmath,amssymb,bm}
\usepackage{booktabs,longtable,array,tabularx}
\usepackage{geometry,xcolor,hyperref}
\usepackage{enumitem}

\geometry{left=2.2cm,right=2.2cm,top=2.2cm,bottom=2.2cm}
\hypersetup{colorlinks=true,linkcolor=blue!55!black,citecolor=blue!55!black,urlcolor=blue!55!black}
\setlength{\parindent}{2em}
\setlength{\parskip}{0.35em}
\renewcommand{\arraystretch}{1.12}

\newcommand{\vect}[1]{\bm{#1}}
\newcommand{\important}[1]{%
  \begin{center}
  \setlength{\fboxsep}{8pt}%
  \fcolorbox{blue!55!black}{blue!3}{\parbox{0.92\textwidth}{#1}}
  \end{center}}

\title{第五题：不可压缩 Navier--Stokes 方程\\自包含题目}
\author{}
\date{}

\begin{document}
\maketitle

\important{\centering\textbf{说明：本文仅给出题目，不包含解答。}原题的控制方程、两种求解方法和两个数值算例保持不变；为使题目能够脱离参考文献独立使用，下文把文献\cite{Ghia1982,Erturk2005,Kohno2006,McQuain1994,Ribbens1994,Jagannathan2014,Erturk2007}中本题验证所需的几何、坐标约定和基准数据直接列入题面。}

\section{控制方程}

考虑二维、等密度、不可压缩 Newton 流体：
\begin{align}
  \nabla\cdot\vect{u} &= 0, \label{eq:continuity}\\
  \frac{\partial\vect{u}}{\partial t}
  +\nabla\cdot(\vect{u}\otimes\vect{u})
  &= -\frac{1}{\rho}\nabla p
  \frac{1}{\rho}\nabla\cdot(\mu\nabla\vect{u}). \label{eq:momentum}
\end{align}
其中，\(\vect{u}=(u,v)^{\mathsf T}\) 为速度，\(p\) 为压力，\(\rho\) 为常密度，\(\mu\) 为动力黏度，\(\nu=\mu/\rho\) 为运动黏度。允许在程序中采用 OpenFOAM 的运动学压力
\[
  p_k=\frac{p}{\rho},
\]
此时动量方程右端写成 \(-\nabla p_k+\nabla\cdot(\nu\nabla\vect{u})\)。

\important{\[
\boxed{\nabla\cdot\vect{u}=0,\qquad
\frac{\partial\vect{u}}{\partial t}
+\nabla\cdot(\vect{u}\otimes\vect{u})
=-\nabla p_k+\nabla\cdot(\nu\nabla\vect{u})}
\]}

请分别基于投影法\cite{Chorin1967,Guermond2006}和 PISO 方法\cite{Issa1986,Jasak1996}写出不可压缩 Navier--Stokes 方程的完整求解过程，并采用有限体积法以及 OpenFOAM 提供的离散格式，各开发一个数值求解器。

\section{统一的有限体积记号}

两种方法的推导、公式和程序说明均须使用下列统一记号：

\begin{center}
\begin{tabularx}{0.94\textwidth}{>{\centering\arraybackslash}p{0.20\textwidth}X}
\toprule
记号 & 含义 \\
\midrule
\(c\), \(V_c\) & 控制体及其体积（二维计算取单位厚度） \\
\(f\) & 内部面 \\
\(O_f,N_f\) & 面 \(f\) 的 owner 单元和 neighbour 单元 \\
\(\vect{n}_f\) & 从 \(O_f\) 指向 \(N_f\) 的全局单位法向量 \\
\(\vect{S}_f=|S_f|\vect{n}_f\) & 按 \(O_f\to N_f\) 定向的全局有向面积向量 \\
\(\sigma_{cf}\) & \(c=O_f\) 时为 \(+1\)，\(c=N_f\) 时为 \(-1\) \\
\(\vect{S}_{cf}=\sigma_{cf}\vect{S}_f\) & 相对于控制体 \(c\) 的外向面积向量 \\
\(\Phi_f=\vect{u}_f\cdot\vect{S}_f\) & 按 \(O_f\to N_f\) 定向的体积通量 \\
\(Q_c^{\mathrm{mass}}\) & 控制体 \(c\) 上尚未除以体积的离散质量通量和 \\
\(R_c^{\mathrm{mass}}=Q_c^{\mathrm{mass}}/V_c\) & 除以体积后的离散质量残差 \\
\bottomrule
\end{tabularx}
\end{center}

内部面通量只计算一次，并按 owner 加、neighbour 减进行装配；因此连续性方程的离散残差须写成
\begin{align}
  Q_c^{\mathrm{mass}}
  &=\sum_{f\in\mathcal F_c^{\mathrm{int}}}\sigma_{cf}\Phi_f
    +Q_c^{\mathrm{bnd}},\\
  R_c^{\mathrm{mass}}&=\frac{Q_c^{\mathrm{mass}}}{V_c}.
\end{align}

\important{\[
\boxed{\vect{S}_f=|S_f|\vect{n}_f,\qquad
\vect{S}_{cf}=\sigma_{cf}\vect{S}_f,\qquad
\Phi_f=\vect{u}_f\cdot\vect{S}_f,\qquad
R_c^{\mathrm{mass}}=\frac{1}{V_c}\sum_f\sigma_{cf}\Phi_f}
\]}

全文不得使用 \(d(f)\) 表示相邻单元；相邻关系一律直接使用 \(O_f,N_f\)。面面积不得另记为 \(A_f\)，而应写为 \(|S_f|\)。

\section{两种方法的任务要求}

\subsection{投影法}

依据 Chorin\cite{Chorin1967}及 Guermond、Minev 和 Shen\cite{Guermond2006}，从式\eqref{eq:continuity}--\eqref{eq:momentum}出发，详细完成下列工作：

\begin{enumerate}[label=\arabic*.]
  \item 给出速度预测方程，明确对流项、扩散项和已知压力项的时间层；
  \item 从速度校正式和 \(\nabla\cdot\vect{u}^{n+1}=0\) 推导压力或压力增量的 Poisson 方程；
  \item 给出单元速度、运动学压力及面通量的校正关系；
  \item 对上述每一步进行有限体积离散，说明时间项、对流通量、扩散通量、压力梯度和非正交修正的处理；
  \item 说明速度边界、压力边界以及压力零空间的处理；
  \item 将上述步骤对应到 OpenFOAM 的场、矩阵、面通量和线性方程求解操作，形成可运行求解器。
\end{enumerate}

所采用的增量投影形式最终应满足下列方法定义；题目要求给出从控制方程到这些关系的完整推导，而不是只抄写结果：
\begin{align}
 \frac{\vect{u}^{*}-\vect{u}^{n}}{\Delta t}
 +\nabla\cdot(\vect{u}\otimes\vect{u})^{\mathrm{disc}}
 &=-\nabla p_k^n+\nabla\cdot(\nu\nabla\vect{u})^{\mathrm{disc}},\\
 \nabla^2\pi^{n+1}&=\frac{1}{\Delta t}\nabla\cdot\vect{u}^{*},\\
 \vect{u}^{n+1}&=\vect{u}^{*}-\Delta t\nabla\pi^{n+1},\\
 p_k^{n+1}&=p_k^n+\pi^{n+1},\\
 \Phi_f^{n+1}&=\Phi_f^{*}-\Delta t\,|S_f|(\nabla_n\pi^{n+1})_f.
\end{align}

\important{\[
\boxed{\nabla^2\pi^{n+1}=\frac{1}{\Delta t}\nabla\cdot\vect{u}^{*},\qquad
\vect{u}^{n+1}=\vect{u}^{*}-\Delta t\nabla\pi^{n+1},\qquad
\nabla\cdot\vect{u}^{n+1}=0}
\]}

\subsection{PISO 方法}

依据 Issa\cite{Issa1986}和 Jasak\cite{Jasak1996}，从离散动量方程出发，在单个时间步内进行多次压力--速度校正。须详细完成下列工作：

\begin{enumerate}[label=\arabic*.]
  \item 将离散动量方程整理为对角系数、非对角贡献和压力梯度的代数形式；
  \item 构造预测速度和预测面通量；
  \item 由离散连续性方程推导 PISO 压力方程；
  \item 给出压力、速度和面通量的第一次校正；
  \item 在同一时间步内至少再进行一次 PISO 校正，并说明各校正量来自哪一层离散动量方程；
  \item 将各步骤对应到 OpenFOAM 的动量矩阵、\(rAU\)、\(HbyA\)、\(phiHbyA\)、压力方程通量及速度修正，形成可运行求解器。
\end{enumerate}

若单元动量方程写成
\[
  A_c\vect{u}_c=\vect{H}_c-V_c(\nabla p_k)_c,
\]
则须从该方程推导并实现与下列 PISO 关系一致的压力--速度耦合：
\begin{align}
  rA_c&=A_c^{-1},\qquad
  \vect{HbyA}_c=rA_c\vect{H}_c,\\
  \vect{u}_c&=\vect{HbyA}_c-rA_cV_c(\nabla p_k)_c,\\
  \nabla\cdot\left(rA\,V\nabla p_k\right)&=\nabla\cdot\vect{HbyA},
\end{align}
并以压力方程的离散面通量对 \(\Phi_f\) 进行守恒校正。

\important{\[
\boxed{A_c\vect{u}_c=\vect{H}_c-V_c(\nabla p_k)_c,\qquad
\nabla\cdot\left(rA\,V\nabla p_k\right)=\nabla\cdot\vect{HbyA},\qquad
\sum_f\sigma_{cf}\Phi_f=0}
\]}

\subsection{OpenFOAM 空间离散与共同检查项}

两个求解器均须采用 OpenFOAM 提供的空间离散操作。至少应明确：

\begin{itemize}
  \item 对流项的面速度插值及有界/迎风处理；
  \item 扩散项的 Gauss 离散与非正交修正；
  \item 压力梯度、速度散度和压力 Laplacian 的离散；
  \item 面通量 \(\Phi_f\) 的构造与压力校正；
  \item 固壁不可穿透条件 \(\Phi_b=0\)，以及压力参考单元和参考值；
  \item 每个时间步结束后 \(R_c^{\mathrm{mass}}\) 的检查。
\end{itemize}

两种求解器必须使用相同的网格、边界条件、时间步控制和空间离散设置完成对照实验；不得把 PISO 简化为重复的 Chorin 投影，也不得用稳态 SIMPLE 外迭代替代题目要求的 PISO 时间步内多重校正。

\section{数值算例一：方形顶盖驱动腔流}

\subsection{几何、参数与边界条件}

计算域为
\[
  \Omega=[0,1]\times[0,1].
\]
上壁 \(y=1\) 以单位速度向右运动，其他三面为静止无滑移壁面：
\[
  \vect{u}=(1,0)^{\mathsf T}\quad(y=1),\qquad
  \vect{u}=\vect{0}\quad(x=0,\ x=1,\ y=0).
\]
取 \(U=L=1\)，故
\[
  Re=\frac{UL}{\nu}=\frac{1}{\nu}.
\]
在 \(Re=1000,3200,5000\) 三种工况下计算，对应
\[
  \nu=10^{-3},\quad 3.125\times10^{-4},\quad 2\times10^{-4}.
\]
各工况均使用不少于三种逐步加密的混合非结构网格，直至稳态。压力采用齐次法向梯度边界条件并给定一个参考压力自由度。初始场可取 \(\vect{u}=\vect{0}\)、\(p_k=0\)。

\subsection{须提交的结果}

对投影法求解器和 PISO 求解器分别给出：

\begin{enumerate}[label=\arabic*.]
  \item 三种 Reynolds 数、各网格分辨率下的稳态流线图；
  \item 主涡及各角涡的中心位置；若计算流函数和涡量，统一采用
  \[
    u=\frac{\partial\psi}{\partial y},\qquad
    v=-\frac{\partial\psi}{\partial x},\qquad
    \omega=\frac{\partial v}{\partial x}-\frac{\partial u}{\partial y}=-\nabla^2\psi;
  \]
  \item 几何中心竖线上的水平速度 \(u(0.5,y)\)；
  \item 几何中心横线上的竖直速度 \(v(x,0.5)\)；
  \item 与文献\cite{Ghia1982,Erturk2005}数据的曲线及表格比较；
  \item 网格加密变化以及离散质量残差 \(R_c^{\mathrm{mass}}\)；
  \item 两个求解器在同一设置下的结果差异。
\end{enumerate}

在流线拓扑上，\(Re=1000\) 应至少分辨主涡、左下角涡和右下角涡；从 \(Re=3200\) 起应能观察到左上角涡\cite{Ghia1982,Erturk2005}。

\subsection{方腔基准数据}

表\ref{tab:ghia-centres}给出 Ghia 等\cite{Ghia1982}的涡心坐标。P 表示主涡，BR1、BL1 和 TL1 分别表示第一右下、左下和左上角涡；``--'' 表示该文献在相应工况下未给出该涡心。

\begin{table}[htbp]
\centering
\caption{Ghia 等\cite{Ghia1982}的方腔涡心坐标 \((x,y)\)}
\label{tab:ghia-centres}
\begin{tabular}{lccc}
\toprule
涡 & \(Re=1000\) & \(Re=3200\) & \(Re=5000\) \\
\midrule
P   & (0.5313, 0.5625) & (0.5165, 0.5469) & (0.5117, 0.5352) \\
BR1 & (0.8594, 0.1094) & (0.8125, 0.0859) & (0.8086, 0.0742) \\
BL1 & (0.0859, 0.0781) & (0.0859, 0.1094) & (0.0703, 0.1367) \\
TL1 & --               & (0.0547, 0.8984) & (0.0625, 0.9102) \\
\bottomrule
\end{tabular}
\end{table}

表\ref{tab:erturk-primary}给出 Erturk、Corke 和 G\"ok\c{c}\"ol\cite{Erturk2005}在 \(601\times601\) 网格上的主涡数据，用于同时校验主涡强度、涡量及位置。

\begin{table}[htbp]
\centering
\caption{Erturk 等\cite{Erturk2005}的方腔主涡数据}
\label{tab:erturk-primary}
\begin{tabular}{cccc}
\toprule
\(Re\) & \(\psi_{\min}\) & \(\omega\) & \((x,y)\) \\
\midrule
1000 & -0.118781 & -2.065530 & (0.5300, 0.5650) \\
5000 & -0.121289 & -1.926601 & (0.5150, 0.5350) \\
\bottomrule
\end{tabular}
\end{table}

\begin{longtable}{r|rrr}
\caption{Ghia 等\cite{Ghia1982}：\(x=0.5\) 上的 \(u(0.5,y)\)}
\label{tab:ghia-u}\\
\toprule
\(y\) & \(Re=1000\) & \(Re=3200\) & \(Re=5000\) \\
\midrule
\endfirsthead
\toprule
\(y\) & \(Re=1000\) & \(Re=3200\) & \(Re=5000\) \\
\midrule
\endhead
1.0000 & 1.00000 & 1.00000 & 1.00000 \\
0.9766 & 0.65928 & 0.53236 & 0.48223 \\
0.9688 & 0.57492 & 0.48296 & 0.46120 \\
0.9609 & 0.51117 & 0.46547 & 0.45992 \\
0.9531 & 0.46604 & 0.46101 & 0.46036 \\
0.8516 & 0.33304 & 0.34682 & 0.33556 \\
0.7344 & 0.18719 & 0.19791 & 0.20087 \\
0.6172 & 0.05702 & 0.07156 & 0.08183 \\
0.5000 & -0.06080 & -0.04272 & -0.03039 \\
0.4531 & -0.10648 & -0.086636\textsuperscript{*} & -0.07404 \\
0.2813 & -0.27805 & -0.24427 & -0.22855 \\
0.1719 & -0.38289 & -0.34323 & -0.33050 \\
0.1016 & -0.29730 & -0.41933 & -0.40435 \\
0.0703 & -0.22220 & -0.37827 & -0.43643 \\
0.0625 & -0.20196 & -0.35344 & -0.42901 \\
0.0547 & -0.18109 & -0.32407 & -0.41165 \\
0.0000 & 0.00000 & 0.00000 & 0.00000 \\
\bottomrule
\multicolumn{4}{p{0.88\textwidth}}{\footnotesize \textsuperscript{*}Ghia 等原表在 \(Re=3200,y=0.4531\) 处印为 \(-0.86636\)。该值与相邻点及速度曲线明显不连续；本题比较时采用通常使用的勘误值 \(-0.086636\)，并要求在报告中说明所用版本。}
\end{longtable}

\begin{longtable}{r|rrr}
\caption{Ghia 等\cite{Ghia1982}：\(y=0.5\) 上的 \(v(x,0.5)\)}
\label{tab:ghia-v}\\
\toprule
\(x\) & \(Re=1000\) & \(Re=3200\) & \(Re=5000\) \\
\midrule
\endfirsthead
\toprule
\(x\) & \(Re=1000\) & \(Re=3200\) & \(Re=5000\) \\
\midrule
\endhead
1.0000 & 0.00000 & 0.00000 & 0.00000 \\
0.9688 & -0.21388 & -0.39017 & -0.49774 \\
0.9609 & -0.27669 & -0.47425 & -0.55069 \\
0.9531 & -0.33714 & -0.52357 & -0.55408 \\
0.9453 & -0.39188 & -0.54053 & -0.52876 \\
0.9063 & -0.51500 & -0.44307 & -0.41442 \\
0.8594 & -0.42665 & -0.37401 & -0.36214 \\
0.8047 & -0.31966 & -0.31184 & -0.30018 \\
0.5000 & 0.02526 & 0.00999 & 0.00945 \\
0.2344 & 0.32235 & 0.28188 & 0.27280 \\
0.2266 & 0.33075 & 0.29030 & 0.28066 \\
0.1563 & 0.37095 & 0.37119 & 0.35368 \\
0.0938 & 0.32627 & 0.42768 & 0.42951 \\
0.0781 & 0.30353 & 0.41906 & 0.43648 \\
0.0703 & 0.29012 & 0.40917 & 0.43329 \\
0.0625 & 0.27485 & 0.39560 & 0.42447 \\
0.0000 & 0.00000 & 0.00000 & 0.00000 \\
\bottomrule
\end{longtable}

为避免只在 Ghia 等\cite{Ghia1982}的采样点上比较，另用 Erturk 等\cite{Erturk2005}的细网格数据进行交叉验证。

\begin{longtable}{r|rr}
\caption{Erturk 等\cite{Erturk2005}：\(x=0.5\) 上的 \(u(0.5,y)\)}
\label{tab:erturk-u}\\
\toprule
\(y\) & \(Re=1000\) & \(Re=5000\) \\
\midrule
\endfirsthead
\toprule
\(y\) & \(Re=1000\) & \(Re=5000\) \\
\midrule
\endhead
1.000 & 1.0000 & 1.0000 \\
0.990 & 0.8486 & 0.6866 \\
0.980 & 0.7065 & 0.5159 \\
0.970 & 0.5917 & 0.4749 \\
0.960 & 0.5102 & 0.4739 \\
0.950 & 0.4582 & 0.4738 \\
0.940 & 0.4276 & 0.4683 \\
0.930 & 0.4101 & 0.4582 \\
0.920 & 0.3993 & 0.4452 \\
0.910 & 0.3913 & 0.4307 \\
0.900 & 0.3838 & 0.4155 \\
0.500 & -0.0620 & -0.0319 \\
0.200 & -0.3756 & -0.3100 \\
0.180 & -0.3869 & -0.3285 \\
0.160 & -0.3854 & -0.3467 \\
0.140 & -0.3690 & -0.3652 \\
0.120 & -0.3381 & -0.3876 \\
0.100 & -0.2960 & -0.4168 \\
0.080 & -0.2472 & -0.4419 \\
0.060 & -0.1951 & -0.4272 \\
0.040 & -0.1392 & -0.3480 \\
0.020 & -0.0757 & -0.2223 \\
0.000 & 0.0000 & 0.0000 \\
\bottomrule
\end{longtable}

\begin{longtable}{r|rr}
\caption{Erturk 等\cite{Erturk2005}：\(y=0.5\) 上的 \(v(x,0.5)\)}
\label{tab:erturk-v}\\
\toprule
\(x\) & \(Re=1000\) & \(Re=5000\) \\
\midrule
\endfirsthead
\toprule
\(x\) & \(Re=1000\) & \(Re=5000\) \\
\midrule
\endhead
1.000 & 0.0000 & 0.0000 \\
0.985 & -0.0973 & -0.2441 \\
0.970 & -0.2173 & -0.5019 \\
0.955 & -0.3400 & -0.5700 \\
0.940 & -0.4417 & -0.5139 \\
0.925 & -0.5052 & -0.4595 \\
0.910 & -0.5263 & -0.4318 \\
0.895 & -0.5132 & -0.4147 \\
0.880 & -0.4803 & -0.3982 \\
0.865 & -0.4407 & -0.3806 \\
0.850 & -0.4028 & -0.3624 \\
0.500 & 0.0258 & 0.0117 \\
0.150 & 0.3756 & 0.3699 \\
0.135 & 0.3705 & 0.3878 \\
0.120 & 0.3605 & 0.4070 \\
0.105 & 0.3460 & 0.4260 \\
0.090 & 0.3273 & 0.4403 \\
0.075 & 0.3041 & 0.4426 \\
0.060 & 0.2746 & 0.4258 \\
0.045 & 0.2349 & 0.3868 \\
0.030 & 0.1792 & 0.3263 \\
0.015 & 0.1019 & 0.2160 \\
0.000 & 0.0000 & 0.0000 \\
\bottomrule
\end{longtable}

\section{数值算例二：等边三角形顶盖驱动腔流}

\subsection{几何、参数与边界条件}

采用原题坐标系，三角形三个顶点为
\[
  A=(-\sqrt{3},0),\qquad B=(\sqrt{3},0),\qquad C=(0,-3).
\]
顶边 \(AB\) 以单位速度向右运动，左右斜边均为静止无滑移壁面：
\[
  \vect{u}=(1,0)^{\mathsf T}\quad(y=0),\qquad
  \vect{u}=\vect{0}\quad\text{（两条斜边）}.
\]
Reynolds 数定义为
\[
  Re=\frac{Ua}{\nu},
\]
其中 \(U=1\)、\(a=1\)。计算 \(Re=100,200,500\)，对应
\[
  \nu=0.01,\qquad 0.005,
  \qquad 0.002.
\]
各工况均在逐步加密的混合非结构网格上求至稳态；其余初始条件、压力参考和两种求解器的对照原则与方腔算例相同。

文献\cite{Erturk2007}使用的等边三角形顶点为
\[
  (-\sqrt3,1),\quad(\sqrt3,1),\quad(0,-2),
\]
仅比本题坐标系整体向上平移 1。因此将文献\cite{Erturk2007}中的坐标换算到本题时须采用
\important{\[
\boxed{x_{\text{题}}=x_{[18]},\qquad y_{\text{题}}=y_{[18]}-1.}
\]}

\subsection{须提交的结果}

对投影法求解器和 PISO 求解器分别给出：

\begin{enumerate}[label=\arabic*.]
  \item \(Re=100,200,500\) 在各网格分辨率下的稳态流线；
  \item 各可分辨涡旋的中心位置；
  \item 主涡的 \(\psi\)、\(\omega\) 及涡心坐标；
  \item 中心线 \(x=0\) 上的水平速度 \(u(0,y)\)；
  \item 横线 \(y=-1\) 上的竖直速度 \(v(x,-1)\)；
  \item 与文献\cite{Kohno2006,McQuain1994,Ribbens1994,Erturk2007}数据的曲线及表格比较；
  \item 网格加密变化、离散质量残差以及两种求解器的差异。
\end{enumerate}

文献\cite{Jagannathan2014}研究的是\emph{直角三角形}顶盖驱动腔流，其几何及取样坐标与本题等边三角形不同；可用于方法和流动现象的定性参考，但不得把其中坐标或中心线速度直接作为本题的定量基准。

\subsection{三角腔基准数据}

Kohno 和 Bathe\cite{Kohno2006}使用的涡旋编号为 \(C_1,C_2,C_3,C_4\)。表\ref{tab:tri-centres}列出本题三种 Reynolds 数下的涡心坐标；``--'' 表示对应涡旋尚未出现或文献未报告。

\begin{table}[htbp]
\centering
\caption{Kohno 和 Bathe\cite{Kohno2006}的等边三角腔涡心坐标 \((x,y)\)}
\label{tab:tri-centres}
\begin{tabular}{lccc}
\toprule
涡 & \(Re=100\) & \(Re=200\) & \(Re=500\) \\
\midrule
\(C_1\) & (0.3467, -0.6416) & (0.2133, -0.7265) & (0.1391, -0.7864) \\
\(C_2\) & (0.0382, -2.1146) & (-0.0275, -2.0001) & (-0.1580, -1.9631) \\
\(C_3\) & (0.0002, -2.8254) & (-0.0002, -2.7991) & (-0.0022, -2.7716) \\
\(C_4\) & -- & -- & (-1.3484, -0.5538) \\
\bottomrule
\end{tabular}
\end{table}

表\ref{tab:tri-primary}给出 Erturk 和 G\"ok\c{c}\"ol\cite{Erturk2007}的细网格主涡数据，以及 McQuain 等和 Ribbens 等\cite{McQuain1994,Ribbens1994}的对照数据。文献\cite{Erturk2007}的涡心坐标已经按 \(y_{\text{题}}=y_{[18]}-1\) 换算到本题坐标系。

\begin{table}[htbp]
\centering
\caption{等边三角腔主涡基准数据}
\label{tab:tri-primary}
\begin{tabular}{clccc}
\toprule
\(Re\) & 来源 & \(\psi\) & \(\omega\) & \((x,y)\) \\
\midrule
100 & 文献\cite{Erturk2007} & -0.2482 & -1.3669 & (0.3315, -0.6445) \\
100 & 文献\cite{McQuain1994,Ribbens1994} & -0.247 & -1.373 & (0.329, -0.645) \\
200 & 文献\cite{Erturk2007} & -0.2624 & -1.2518 & (0.2030, -0.7266) \\
200 & 文献\cite{McQuain1994,Ribbens1994} & -0.260 & -1.272 & (0.208, -0.720) \\
500 & 文献\cite{Erturk2007} & -0.2774 & -1.1791 & (0.1319, -0.7793) \\
500 & 文献\cite{McQuain1994,Ribbens1994} & -0.269 & -1.250 & (0.173, -0.735) \\
\bottomrule
\end{tabular}
\end{table}

\begin{longtable}{r|rrr}
\caption{Kohno 和 Bathe\cite{Kohno2006}：中心线 \(x=0\) 上的 \(u(0,y)\)}
\label{tab:tri-u}\\
\toprule
\(y\) & \(Re=100\) & \(Re=200\) & \(Re=500\) \\
\midrule
\endfirsthead
\toprule
\(y\) & \(Re=100\) & \(Re=200\) & \(Re=500\) \\
\midrule
\endhead
0.0000  & 1.00000  & 1.00000  & 1.00000 \\
-0.0571 & 0.81766  & 0.78999  & 0.73838 \\
-0.0714 & 0.77171  & 0.73771  & 0.67965 \\
-0.0857 & 0.72676  & 0.68813  & 0.62977 \\
-0.1000 & 0.68320  & 0.64201  & 0.58927 \\
-0.4429 & 0.16395  & 0.23983  & 0.28014 \\
-0.8714 & -0.16653 & -0.10758 & -0.05682 \\
-1.1857 & -0.29476 & -0.33718 & -0.26938 \\
-1.3143 & -0.27349 & -0.37775 & -0.35935 \\
-1.5000 & -0.18897 & -0.29791 & -0.45972 \\
-1.5857 & -0.14445 & -0.22425 & -0.41095 \\
-1.6429 & -0.11669 & -0.17469 & -0.33494 \\
-1.6714 & -0.10376 & -0.15138 & -0.28993 \\
-1.6857 & -0.09758 & -0.14028 & -0.26698 \\
-1.9429 & -0.02046 & -0.01231 & -0.01293 \\
-1.9571 & -0.01798 & -0.00880 & -0.00764 \\
-2.1143 & -0.00032 & 0.01304 & 0.02439 \\
-2.1857 & 0.00312 & 0.01543 & 0.02737 \\
-2.2000 & 0.00355 & 0.01553 & 0.02732 \\
-2.3000 & 0.00481 & 0.01371 & 0.02277 \\
-3.0000 & 0.00000 & 0.00000 & 0.00000 \\
\bottomrule
\end{longtable}

\begin{longtable}{r|rrr}
\caption{Kohno 和 Bathe\cite{Kohno2006}：横线 \(y=-1\) 上的 \(v(x,-1)\)}
\label{tab:tri-v}\\
\toprule
\(x\) & \(Re=100\) & \(Re=200\) & \(Re=500\) \\
\midrule
\endfirsthead
\toprule
\(x\) & \(Re=100\) & \(Re=200\) & \(Re=500\) \\
\midrule
\endhead
1.1547  & 0.00000 & 0.00000 & 0.00000 \\
1.0854  & -0.14663 & -0.20538 & -0.33354 \\
1.0623  & -0.18510 & -0.25872 & -0.40554 \\
1.0277  & -0.23363 & -0.32360 & -0.47331 \\
0.9815  & -0.28188 & -0.38138 & -0.49887 \\
0.9122  & -0.32187 & -0.41260 & -0.45508 \\
0.8545  & -0.33011 & -0.40093 & -0.39707 \\
0.0000  & 0.10458 & 0.09033 & 0.05136 \\
-0.4734 & 0.18995 & 0.26513 & 0.28742 \\
-0.4965 & 0.18981 & 0.26566 & 0.29968 \\
-0.6813 & 0.17807 & 0.23822 & 0.35731 \\
-0.7852 & 0.16183 & 0.20468 & 0.32956 \\
-0.8545 & 0.14537 & 0.17678 & 0.28510 \\
-0.8891 & 0.13501 & 0.16088 & 0.25700 \\
-0.9122 & 0.12723 & 0.14946 & 0.23645 \\
-1.1547 & 0.00000 & 0.00000 & 0.00000 \\
\bottomrule
\end{longtable}

\section{最终提交内容}

须提交一份完整成果，至少包括：

\begin{enumerate}[label=\arabic*.]
  \item 从式\eqref{eq:continuity}--\eqref{eq:momentum}到有限体积代数方程的详细推导；
  \item 投影法完整求解流程、压力 Poisson 方程、速度和面通量校正；
  \item PISO 完整求解流程、动量预测及不少于两次压力--速度校正；
  \item 两个可编译、可运行的 OpenFOAM 有限体积求解器及编译文件；
  \item 方腔和等边三角腔的全部算例文件、网格及关键离散设置；
  \item 两个算例的流线、中心线速度、涡心及必要的 \(\psi,\omega\) 数据；
  \item 与本题所列基准数据的定量对比、网格加密结果和离散质量守恒检查；
  \item 投影法与 PISO 法在相同数值设置下的结果对比。
\end{enumerate}

\important{\centering\textbf{最重要的验收条件：}
两个求解器均须在每个稳态算例中满足离散连续性，正确给出方腔和等边三角腔的流线拓扑，并使中心线速度与涡旋特征随网格加密趋近本题所列文献基准。}

\begin{thebibliography}{99}

\bibitem[8]{Chorin1967}
A. J. Chorin,
``A numerical method for solving incompressible viscous flow problems,''
\emph{Journal of Computational Physics}, 2 (1967), 12--26.

\bibitem[9]{Guermond2006}
J.-L. Guermond, P. Minev and J. Shen,
``An overview of projection methods for incompressible flows,''
\emph{Computer Methods in Applied Mechanics and Engineering}, 195 (2006), 6011--6045.

\bibitem[10]{Issa1986}
R. I. Issa,
``Solution of the implicitly discretised fluid flow equations by operator-splitting,''
\emph{Journal of Computational Physics}, 62 (1986), 40--65.

\bibitem[11]{Jasak1996}
H. Jasak,
\emph{Error Analysis and Estimation for the Finite Volume Method with Applications to Fluid Flows},
PhD thesis, Imperial College London, 1996.

\bibitem[12]{Ghia1982}
U. Ghia, K. N. Ghia and C. T. Shin,
``High-Re solutions for incompressible flow using the Navier--Stokes equations and a multigrid method,''
\emph{Journal of Computational Physics}, 48 (1982), 387--411.

\bibitem[13]{Erturk2005}
E. Erturk, T. C. Corke and C. G\"ok\c{c}\"ol,
``Numerical solutions of 2-D steady incompressible driven cavity flow at high Reynolds numbers,''
\emph{International Journal for Numerical Methods in Fluids}, 48 (2005), 747--774.

\bibitem[14]{Kohno2006}
H. Kohno and K. J. Bathe,
``A flow-condition-based interpolation finite element procedure for triangular grids,''
\emph{International Journal for Numerical Methods in Fluids}, 51 (2006), 673--699.

\bibitem[15]{McQuain1994}
W. D. McQuain, C. J. Ribbens, C.-Y. Wang and L. T. Watson,
``Steady viscous flow in a trapezoidal cavity,''
\emph{Computers \& Fluids}, 23 (1994), 613--626.

\bibitem[16]{Ribbens1994}
C. J. Ribbens, L. T. Watson and C.-Y. Wang,
``Steady viscous flow in a triangular cavity,''
\emph{Journal of Computational Physics}, 112 (1994), 173--181.

\bibitem[17]{Jagannathan2014}
A. Jagannathan, R. Mohan and M. Dhanak,
``A spectral method for the triangular cavity flow,''
\emph{Computers \& Fluids}, 95 (2014), 40--48.

\bibitem[18]{Erturk2007}
E. Erturk and O. Gokcol,
``Fine grid numerical solutions of triangular cavity flow,''
\emph{European Physical Journal Applied Physics}, 38 (2007), 97--105.

\end{thebibliography}

\end{document}
